\documentclass{article}
\usepackage[utf8]{inputenc}
\usepackage{geometry}
\usepackage{hyperref}
\usepackage{minted}

\geometry{a4paper, margin=1in}

\title{Password Strength Analyzer: \\ \large Explained for 5th Graders!}
\author{Machine Learning Project}
\date{\today}

\begin{document}

\maketitle

\section*{Introduction}
Hey there! Have you ever wondered how a website knows if your newly typed password is "weak" or "strong"? We wrote a computer program to do exactly that! 

We are using something called \textbf{Machine Learning}, which means we teach the computer by showing it millions of examples instead of just giving it a list of rules. Let's look at what each part of our code does!

\section*{Step 1: Getting Our Tools Ready}
Just like you need a pencil or a calculator for math, our computer program needs special tools.

\begin{itemize}
    \item \texttt{import pandas as pd} - We bring in a tool named "pandas" (not the fluffy animal!). It helps us read and organize big spreadsheets of data.
    \item \texttt{import numpy as np} - "numpy" is a tool that helps the computer do super-fast math with lists of numbers.
    \item \texttt{from sklearn... import TfidfVectorizer} - Computers don't understand letters; they only understand numbers! This tool safely changes text (like passwords) into numbers.
    \item \texttt{from sklearn... import train\_test\_split} - Think of studying for a test. You use some flashcards to study (train), and the teacher uses other questions to quiz you (test). This tool splits our data into "study" flashcards and "quiz" questions.
    \item \texttt{from sklearn... import RandomForestClassifier} - This is our computer's "brain"! It works by asking lots of yes/no questions to guess the answer.
    \item \texttt{import getpass} - This tool hides your password with stars (***) when you type it so nobody can spy on it!
\end{itemize}

\section*{Step 2: Reading the Data}
\begin{itemize}
    \item \texttt{data = pd.read\_csv("data.csv", on\_bad\_lines='skip')} - We tell "pandas" to open our giant file (\texttt{data.csv}) that has tons of passwords. If it finds a broken line, we tell it to skip it.
    \item \texttt{print(data.head())} - This simply shows us the first 5 lines of our spreadsheet so we know it loaded correctly.
\end{itemize}

\section*{Step 3: Cleaning Up the Data}
\begin{itemize}
    \item \texttt{data = data.dropna()} - Sometimes, our spreadsheet has empty, blank spaces. This line throws away the blank spaces so our computer won't get confused!
    \item \texttt{data["strength"] = data["strength"].map(\{0: "Weak", 1: "Medium", 2: "Strong"\})} - The spreadsheet uses numbers to rate passwords: 0, 1, and 2. We change these numbers into words so they are much easier for humans to read.
    \item \texttt{print(data.sample(5))} - We pick 5 random passwords from our list and display them on the screen.
\end{itemize}

\section*{Step 4: Breaking Passwords into Tiny Pieces}
\begin{minted}{python}
def word(password):
    character=[]
    for i in password:
        character.append(i)
    return character
\end{minted}

This creates a small set of instructions named \texttt{word()}. It takes a password (like "apple123") and breaks it apart into single letters and numbers (like 'a', 'p', 'p', 'l', 'e', '1', '2', '3'). This makes it easier for the computer to search for patterns!

\section*{Step 5: Preparing for the Exam!}
\begin{itemize}
    \item \texttt{x = np.array(data["password"])} - We put all of the passwords into a group called \texttt{x}.
    \item \texttt{y = np.array(data["strength"])} - We put all of the answers (whether they are Weak, Medium, or Strong) into a group called \texttt{y}.
    \item \texttt{tdif = TfidfVectorizer(...)} and \texttt{x = tdif.fit\_transform(x)} - Here, we change every single letter in our \texttt{x} group into math numbers.
    \item \texttt{xtrain, xtest, ytrain, ytest = train\_test\_split(...)} - The computer splits everything. 95\% of the passwords will be used to learn and study (\texttt{train}), and the remaining 5\% will be hidden away and used as the final quiz (\texttt{test}).
\end{itemize}

\section*{Step 6: Brain Training!}
\begin{itemize}
    \item \texttt{print("Training...")} - Just telling you to wait!
    \item \texttt{model = RandomForestClassifier(n\_estimators=20, n\_jobs=-1, random\_state=42)} - We build the computer's brain. \texttt{n\_estimators=20} tells it to grow a "forest" of 20 decision-making trees. \texttt{n\_jobs=-1} tells the computer to use all its energy to work as fast as possible.
    \item \texttt{model.fit(xtrain, ytrain)} - \textbf{THIS IS THE MAGIC LINE!} This tells the computer: "Look at the study flashcards (\texttt{xtrain}), look at the answers (\texttt{ytrain}), and figure out the patterns!" The computer is literally learning right now.
    \item \texttt{print("Model Accuracy:", model.score(xtest, ytest))} - The learning is over! Now we give the computer that final 5\% quiz it has never seen before and print out its test score.
\end{itemize}

\section*{Step 7: Do It Yourself!}
\begin{itemize}
    \item \texttt{user = getpass.getpass("Enter Password to test: ")} - Puts up a box for YOU to type a secret password.
    \item \texttt{user\_data = tdif.transform([user]).toarray()} - Converts the word you typed into mathematical numbers.
    \item \texttt{output = model.predict(user\_data)} - The computer uses its trained brain to predict whether your password is Weak, Medium, or Strong.
    \item \texttt{print("Predicted Strength:", output[0])} - Finally, the computer prints out its guess for your password!
\end{itemize}

\end{document}
